\section{Лекция 1. Инструментарий разработчика}

\subsection*{Синопсис}
\paragraph*{Цель}

Настроить рабочее место, понять, что Python-код живёт в экосистеме.
\paragraph*{Длительность}

3–5 дней.
\paragraph*{Теория (минимум)}
\begin{itemize}
\item Что такое интерпретатор, REPL, байт-код, \_\_pycache\_\_.
\item Версии Python: 3.10+, почему не 2.x.
\item python, python3, pip, pipx, uv, poetry.
\item venv: зачем изоляция зависимостей.
\item pyproject.toml vs requirements.txt vs setup.py.
\item Git, .gitignore для Python (.venv, \verb| __pycache__|,  \verb|.pytest_cache|).
\item Редактор: VS Code + Python extension или PyCharm.
\item Terminal: bash/zsh, базовые командыю.
\end{itemize}

\subsection{Что такое интерпретатор, REPL, байт-код , \_\_pycache\_\_ }
\subsubsection{Интерпретатор}

\paragraph*{Python - интерпретируемый язык}
Это значит, что исходный код \texttt{.py} \textbf{не компилируется} в машинный код (python3) напрямую,  а \textbf{выполняется} специальной программой --- \linebreak интерпретатором.
Но это упрощение. На самом деле путь такой: (см. Рис.1).
То есть Python сначала компилирует исходник в байт-код, а потом интерпретирует байт-код виртуальной машиной.
Поэтому Python иногда называют \textbf{«компилируемо-интерпретируемым»} языком.

Что нужно понимать:
\begin{itemize}
\item \textbf{Интерпретатор} — это программа \texttt{python3}, которая читает \texttt{.py}, компилирует в байт-код и выполняет.
\item \textbf{Интерпретатор} не создаёт \texttt{.exe}/бинарник. Чтобы запустить программу, нужен \linebreak  установленный Python.
\item \textbf{Python-код} \underline{медленнее} C, потому что каждая операция проходит через PVM.
\end{itemize}

\begin{figure}[htbp]
  \centering
  \resizebox{0.65\textwidth}{!}{% generated by Plantuml 1.2026.8       
\definecolor{plantucolor0000}{RGB}{0,0,0}%
\definecolor{plantucolor0001}{RGB}{34,34,34}%
\definecolor{plantucolor0002}{RGB}{241,241,241}%
\definecolor{plantucolor0003}{RGB}{24,24,24}%
\scalebox{3.125}{%
\begin{tikzpicture}[yscale=-1
,pstyle0/.style={color=plantucolor0001,fill=plantucolor0001,line width=1.0pt}
,pstyle1/.style={color=plantucolor0003,fill=plantucolor0002,line width=0.5pt}
,pstyle3/.style={color=plantucolor0003,line width=1.0pt}
,pstyle4/.style={color=plantucolor0003,fill=plantucolor0003,line width=1.0pt}
]
\node at (10pt,10pt)[below right,color=black,inner sep=0]{\textbf{Схема обработки и выполнения программы на Python}};
\draw[pstyle0] (174.8pt,51pt) ellipse (10pt and 10pt);
\draw[pstyle1] (119.3pt,123pt) arc (180:270:22pt) -- (141.3pt,101pt) -- (208.35pt,101pt) arc (270:360:22pt) -- (230.35pt,123pt) -- (230.35pt,123pt) arc (0:90:22pt) -- (208.35pt,145pt) -- (141.3pt,145pt) arc (90:180:22pt) -- (119.3pt,123pt) -- cycle;
\node at (129.3pt,111pt)[below right,color=black,inner sep=0]{source.py};
\node at (129.3pt,123pt)[below right,color=black,inner sep=0]{(исходный текст)};
\draw[pstyle1] (122.8pt,201pt) arc (180:270:16pt) -- (138.8pt,185pt) -- (210.575pt,185pt) arc (270:360:16pt) -- (226.575pt,201pt) -- (226.575pt,201pt) arc (0:90:16pt) -- (210.575pt,217pt) -- (138.8pt,217pt) arc (90:180:16pt) -- (122.8pt,201pt) -- cycle;
\node at (132.8pt,195pt)[below right,color=black,inner sep=0]{Лексер / Парсер};
\draw[pstyle1] (62.8pt,279pt) arc (180:270:22pt) -- (84.8pt,257pt) -- (264.725pt,257pt) arc (270:360:22pt) -- (286.725pt,279pt) -- (286.725pt,279pt) arc (0:90:22pt) -- (264.725pt,301pt) -- (84.8pt,301pt) arc (90:180:22pt) -- (62.8pt,279pt) -- cycle;
\node at (72.8pt,267pt)[below right,color=black,inner sep=0]{AST};
\node at (72.8pt,279pt)[below right,color=black,inner sep=0]{(абстрактное синтаксическое дерево)};
\draw[pstyle1] (131.3pt,357pt) arc (180:270:16pt) -- (147.3pt,341pt) -- (202.725pt,341pt) arc (270:360:16pt) -- (218.725pt,357pt) -- (218.725pt,357pt) arc (0:90:16pt) -- (202.725pt,373pt) -- (147.3pt,373pt) arc (90:180:16pt) -- (131.3pt,357pt) -- cycle;
\node at (141.3pt,351pt)[below right,color=black,inner sep=0]{Компилятор};
\draw[pstyle1] (139.8pt,429pt) arc (180:270:16pt) -- (155.8pt,413pt) -- (193.375pt,413pt) arc (270:360:16pt) -- (209.375pt,429pt) -- (209.375pt,429pt) arc (0:90:16pt) -- (193.375pt,445pt) -- (155.8pt,445pt) arc (90:180:16pt) -- (139.8pt,429pt) -- cycle;
\node at (149.8pt,423pt)[below right,color=black,inner sep=0]{Байт-код};
\draw[pstyle1] (87.8pt,507pt) arc (180:270:22pt) -- (109.8pt,485pt) -- (240.15pt,485pt) arc (270:360:22pt) -- (262.15pt,507pt) -- (262.15pt,507pt) arc (0:90:22pt) -- (240.15pt,529pt) -- (109.8pt,529pt) arc (90:180:22pt) -- (87.8pt,507pt) -- cycle;
\node at (97.8pt,495pt)[below right,color=black,inner sep=0]{Виртуальная машина Python};
\node at (97.8pt,507pt)[below right,color=black,inner sep=0]{(PVM)};
\draw[pstyle1] (98.3pt,585pt) arc (180:270:16pt) -- (114.3pt,569pt) -- (235.05pt,569pt) arc (270:360:16pt) -- (251.05pt,585pt) -- (251.05pt,585pt) arc (0:90:16pt) -- (235.05pt,601pt) -- (114.3pt,601pt) arc (90:180:16pt) -- (98.3pt,585pt) -- cycle;
\node at (108.3pt,579pt)[below right,color=black,inner sep=0]{Выполнение программы};
\draw[color=plantucolor0001,line width=1.5pt] (174.8pt,652pt) ellipse (11pt and 11pt);
\draw[pstyle0] (174.8pt,652pt) ellipse (6pt and 6pt);
\draw[pstyle3] (174.8pt,61.15pt) ..controls (174.8pt,71.16pt) and (174.8pt,82.53pt) .. (174.8pt,95.78pt);
\draw[pstyle4] (174.8pt,100.78pt) -- (178.8pt,91.78pt) -- (174.8pt,95.78pt) -- (170.8pt,91.78pt) -- (174.8pt,100.78pt) -- cycle;
\draw[pstyle3] (174.8pt,145.17pt) ..controls (174.8pt,157.82pt) and (174.8pt,168.59pt) .. (174.8pt,179.96pt);
\draw[pstyle4] (174.8pt,184.96pt) -- (178.8pt,175.96pt) -- (174.8pt,179.96pt) -- (170.8pt,175.96pt) -- (174.8pt,184.96pt) -- cycle;
\draw[pstyle3] (174.8pt,217.02pt) ..controls (174.8pt,228.44pt) and (174.8pt,239.29pt) .. (174.8pt,251.97pt);
\draw[pstyle4] (174.8pt,256.97pt) -- (178.8pt,247.97pt) -- (174.8pt,251.97pt) -- (170.8pt,247.97pt) -- (174.8pt,256.97pt) -- cycle;
\draw[pstyle3] (174.8pt,301.17pt) ..controls (174.8pt,313.82pt) and (174.8pt,324.59pt) .. (174.8pt,335.96pt);
\draw[pstyle4] (174.8pt,340.96pt) -- (178.8pt,331.96pt) -- (174.8pt,335.96pt) -- (170.8pt,331.96pt) -- (174.8pt,340.96pt) -- cycle;
\draw[pstyle3] (174.8pt,373.19pt) ..controls (174.8pt,384.93pt) and (174.8pt,396.11pt) .. (174.8pt,407.84pt);
\draw[pstyle4] (174.8pt,412.84pt) -- (178.8pt,403.84pt) -- (174.8pt,407.84pt) -- (170.8pt,403.84pt) -- (174.8pt,412.84pt) -- cycle;
\draw[pstyle3] (174.8pt,445.02pt) ..controls (174.8pt,456.44pt) and (174.8pt,467.29pt) .. (174.8pt,479.97pt);
\draw[pstyle4] (174.8pt,484.97pt) -- (178.8pt,475.97pt) -- (174.8pt,479.97pt) -- (170.8pt,475.97pt) -- (174.8pt,484.97pt) -- cycle;
\draw[pstyle3] (174.8pt,529.17pt) ..controls (174.8pt,541.82pt) and (174.8pt,552.59pt) .. (174.8pt,563.96pt);
\draw[pstyle4] (174.8pt,568.96pt) -- (178.8pt,559.96pt) -- (174.8pt,563.96pt) -- (170.8pt,559.96pt) -- (174.8pt,568.96pt) -- cycle;
\draw[pstyle3] (174.8pt,601.08pt) ..controls (174.8pt,613.23pt) and (174.8pt,624.99pt) .. (174.8pt,635.66pt);
\draw[pstyle4] (174.8pt,640.66pt) -- (178.8pt,631.66pt) -- (174.8pt,635.66pt) -- (170.8pt,631.66pt) -- (174.8pt,640.66pt) -- cycle;
\end{tikzpicture}%
}%
}
  \caption{Схема обработки и выполнения программы на Python}
  \label{fig:python-pipeline}
\end{figure}

\begin{warningbox2}
\texttt{Интерпретатор} --- Исходный код \texttt{.py} компилируется в байт-код.
\end{warningbox2}

\subsubsection{REPL}

\textbf{REPL = Read-Eval-Print Loop}. Это интерактивный режим интерпретатора.

\begin{lstlisting}
$ python3
Python 3.12.0 (main, ...)
>>> 2 + 2
4
>>> name = "world"
>>> print(f"hello, {name}")
hello, world
>>> exit()
\end{lstlisting}

\textbf{Цикл:}
\begin{itemize}
\item Read — читает строку.
\item Eval — вычисляет.
\item Print — печатает результат.
\item Loop — возвращается к началу.
\end{itemize}

\textbf{Зачем нужен:}
\begin{itemize}
\item Быстро проверить идею.
\item Изучить поведение функции/метода.
\item Отладка маленьких фрагментов.
\item Прототипирование.
\end{itemize}

\textbf{Улучшенные REPL:}
\begin{itemize}
\item ipython — подсветка, автодополнение, ?, ??, \%timeit.
\item bpython, ptpython.
\item Jupyter Notebook — REPL в браузере с сохранением состояния и визуализацией.
\end{itemize}

\textbf{Что нужно понимать:}
\begin{itemize}
\item \textbf{REPL} — не для продакшена, а для экспериментов.
\item В \textbf{REPL} переменные живут до выхода.
\item  \textbf{\_} хранит последний результат.
\end{itemize}

\subsubsection{Байт-код}

\textbf{Байт-код} — промежуточное представление между исходником и машинным кодом. Это инструкции для виртуальной машины Python.

Посмотреть байт-код:
\begin{lstlisting}
import dis

def add(a, b):
    return a + b

dis.dis(add)
\end{lstlisting}

Вывод:
\begin{lstlisting}
  2           0 RESUME                   0
  3           2 LOAD_FAST                0 (a)
              4 LOAD_FAST                1 (b)
              6 BINARY_OP                0 (+)
             10 RETURN_VALUE
\end{lstlisting}

\textbf{Что нужно понимать}:
\begin{itemize}
\item Байт-код — платформо-независимый (один и тот же на Linux/Windows/macOS).
\item Его можно дизассемблировать через dis.
\item Байт-код не является машинным кодом и не переносится между версиями Python (.pyc для 3.11 не подойдёт для 3.12).
\end{itemize}

\subsubsection{\_\_pycache\_\_ и .pyc}

Когда Python импортирует модуль, он кэширует скомпилированный байт-код в файлы .pyc, чтобы не компилировать заново при следующем запуске
\begin{lstlisting}
.project/
├── module.py
└── \_\_pycache\_\_/
    └── module.cpython-312.pyc
\end{lstlisting}

Имя \texttt{.pyc: module.<interpreter>-<version>.pyc}, например \texttt{ module.cpython-312.pyc}.

Когда создаётся:
\begin{itemize}
\item При import module.
\item Не создаётся для скрипта, запущенного как python script.py (главный модуль не кэшируется).
\item Можно отключить: python -B script.py или PYTHONDONTWRITEBYTECODE=1.
\end{itemize}

Что нужно понимать:
\begin{itemize}
\item \texttt{ \_\_pycache\_\_} — артефакт, его не коммитят в git.
\item \texttt{.pyc} ускоряет повторные импорты.
\item Если исходник изменился, .pyc перекомпилируется.
\item Удалять \texttt{ \_\_pycache\_\_} безопасно.
\end{itemize}

\subsubsection{ Что проверить руками}
\begin{lstlisting}
# REPL
python3
>>> import this          # Zen of Python
>>> import dis; dis.dis(lambda x: x+1)
>>> exit()

# Байт-код
python3 -c "import dis; dis.dis(compile('a+b', '<s>', 'eval'))"

# __pycache__
mkdir /tmp/pycache-test && cd /tmp/pycache-test
echo "x = 42" > mod.py
python3 -c "import mod; print(mod.x)"
ls -la __pycache__/
\end{lstlisting}

\subsection{ Версии Python: 3.10+, почему не 2.x2.1.}
 
\subsubsection{Python 2 vs Python 3}

Python 2 вышел в 2000, Python 3 — в 2008. Это несовместимые языки. Python 2 официально мёртв с 1 января 2020 (EOL — End of Life).

\textbf{Ключевые различия:}
\begin{table}[h]
\centering
\begin{tabular}{|l|l|}
\hline
\textbf{Версия} & \textbf{Особенность} \\
\hline
3.10 & match/case, \texttt{X | Y} в типах \\
3.11 & Быстрее на 10--60\% \\
3.12 & Ещё быстрее, новый синтаксис \texttt{type} \\
\hline
\end{tabular}
\caption{Актуальные версии Python}
\label{tab:python-versions}
\end{table}

АспектPython 2Python 3printоператор print "hi"функция print("hi")Строкиstr (bytes) + unicodestr (unicode) + bytesДеление5/2 == 25/2 == 2.5, 5//2 == 2rangeсписокленивый объектinputeval вводастрокаUnicodeпо умолчанию нетпо умолчанию даПоддержкаEOL 2020активная

Почему не 2.x:
Нет обновлений безопасности.
Большинство библиотек (numpy, pandas, fastapi, requests) уже не поддерживают 2.x.
Синтаксис 3 чище.
Python 2 несовместим с современным тулингом.

2.2. Какие версии актуальны

3.10 — минимум для большинства современных проектов (match/case, X | Y в типах).
3.11 — быстрее (~10–60%), TaskGroup, tomllib, лучшие сообщения об ошибках.
3.12 — ещё быстрее, type statement, улучшенный f-string.
3.13 — экспериментальный JIT, улучшенный REPL, free-threaded (без GIL) — пока не для продакшена.

Рекомендация: целевая версия — 3.11 или 3.12. Минимум — 3.10.

2.3. Что нужно понимать
Указывай requires-python = ">=3.11" в pyproject.toml.
Используй python3.12 явно, если нужно.
Не пиши код, который работает только на одной версии, без причины.
Проверяй версию: python3 --version.

2.4. Что проверить руками

python3 --version
python3 -c "import sys; print(sys.version_info)"
# Проверка новых фич

python3 -c "match 1:
    case 1: print('one')"

3. python, python3, pip, pipx, uv, poetry
3.1. python vs python3

python3 — гарантированно Python 3 на Linux/macOS.
python — может быть Python 2, а может не быть вообще.
На Windows python — обычно Python 3 (если установлен).
Правило: всегда используй python3 в скриптах и документации для Linux/macOS.

Проверка:

which python3
which python
python3 --version

3.2. pip

pip = Pip Installs Packages. Менеджер пакетов Python.

pip install requests
pip install "requests>=2.31,<3"
pip install -e .                    # editable install
pip install -r requirements.txt
pip list
pip show requests
pip uninstall requests
pip freeze > requirements.txt

Важно: pip — это не «системный» установщик. Он ставит пакеты в текущее окружение. Без venv — в глобальное, что плохо.

Правило: всегда активируй venv перед pip install.

Альтернативы:

python3 -m pip — гарантирует, что pip от нужного Python.
pip3 — явный Python 3.

3.3. pipxpipx = pip + изоляция. Устанавливает приложения (CLI-утилиты) в изолированные окружения и добавляет их в PATH.

pipx install ruff
pipx install poetry
pipx install black
pipx list
pipx upgrade ruff
pipx uninstall ruff

Когда использовать:
Устанавливаешь CLI-инструмент (ruff, black, poetry, httpie).
Не хочешь засорять глобальный site-packages.
Нужны разные версии инструмента в разных проектах.

Когда не использовать:
Устанавливаешь библиотеку для проекта — тогда pip + venv.

3.4. uvuv — сверхбыстрый менеджер пакетов и окружений от Astral (создатели ruff). Написан на Rust.

# Установка
curl -LsSf https://astral.sh/uv/install.sh | sh

# Создание venv
uv venv
source .venv/bin/activate

# Установка
uv pip install requests
uv pip install -e ".[dev]"
uv pip compile requirements.in -o requirements.txt

# Управление проектом (аналог poetry)
uv init myproject
cd myproject
uv add requests
uv add --dev pytest
uv sync
uv run python main.py
uv lock

Преимущества:
В 10–100 раз быстрее pip.Единый инструмент: venv + pip + lock + resolver.
Совместим с pyproject.toml.
Кэширует агрессивно.

Рекомендация: для новых проектов — uv. Для legacy — pip.

3.5. poetryPoetry — менеджер зависимостей и упаковки. Более старый, чем uv, но зрелый.

pipx install poetry
poetry new myproject
cd myproject
poetry add requests
poetry add --group dev pytest
poetry install
poetry run python main.py
poetry build
poetry publish

Плюсы:

Единый pyproject.toml + poetry.lock.
Хорошая работа с версиями.
Упаковка и публикация.

Минусы:
Медленнее uv.
Свой формат секций в pyproject.toml.

3.6. СравнениеИнструментНазначениеСкоростьКогда использоватьpipустановка пакетовсредняяlegacy, простые случаиpipxCLI-приложениясредняяruff, black, poetryuvвсёочень высокаяновые проектыpoetryзависимости + упаковкасредняяlegacy, публикация

3.7. Что проверить руками

python3 --version
python3 -m pip --version
pipx --version
uv --version
poetry --version

# Создать проект через uv
uv init /tmp/uvtest && cd /tmp/uvtest
uv add requests
uv run python -c "import requests; print(requests.__version__)"

4. venv: зачем изоляция зависимостей
4.1. Проблема
Представь: у тебя два проекта.Проект A требует requests==2.28.Проект B требует requests==2.31.Если ставить глобально — конфликт. Один проект сломается.
Решение: виртуальное окружение — изолированная папка с собственным site-packages.

4.2. Как работает
python3 -m venv .venv
source .venv/bin/activate       # Linux/macOS
# или
.venv\Scripts\activate           # Windows

# Теперь python и pip указывают на .venv
which python                     # /path/to/.venv/bin/python
pip install requests             # ставится в .venv/lib/.../site-packages

deactivate                       # выйтиЧто создаётся:.venv/
├── bin/                  # python, pip, activate
├── lib/
│   └── python3.12/
│       └── site-packages/   # сюда ставятся пакеты
└── pyvenv.cfg

Ключевое: activate меняет PATH и VIRTUAL_ENV, чтобы python и pip указывали на .venv.

4.3. Правила
Один venv на проект.
Не коммитить .venv — добавить в .gitignore.
Не создавать venv внутри venv.
Удалять venv безопасно — просто rm -rf .venv.
Фиксировать зависимости — requirements.txt или uv.lock/poetry.lock.

4.4. Альтернативы
virtualenv — старее, быстрее для Python 2.
conda — для data science, управляет и не-Python зависимостями.
uv venv — быстрее.
poetry shell — активирует venv проекта.

4.5. Что проверить руками

mkdir /tmp/venvtest && cd /tmp/venvtest
python3 -m venv .venv
source .venv/bin/activate
which python
pip install requests
pip list
python -c "import requests; print(requests.__file__)"
deactivate

5. pyproject.toml vs requirements.txt vs setup.py
5.1. setup.py (legacy)
Исторический способ. Исполняемый Python-файл.

from setuptools import setup, find_packages

setup(
    name="myproject",
    version="0.1.0",
    packages=find_packages(),
    install_requires=["requests>=2.31"],
)

Проблемы:
Исполняемый код при установке — небезопасно.
Нельзя прочитать метаданные без запуска.
Дублирование с setup.cfg.

Статус: deprecated, но ещё встречается в legacy-проектах.

5.2. requirements.txt

Простой список зависимостей.

requests==2.31.0
httpx>=0.27,<1
pydantic~=2.6

Плюсы:
Просто.
Понятно.
Работает везде.

Минусы:
Нет метаданных проекта (имя, версия, entry points).
Нет разделения на prod/dev (обычно делают requirements-dev.txt).
Нет lock-файла (нужен pip-tools или pip freeze).

Когда использовать:
Простые проекты.
Legacy.
Docker-образы (иногда).

5.3. pyproject.toml (современный стандарт)

PEP 518 / PEP 621.
Единый файл конфигурации проекта.

[build-system]
requires = ["setuptools>=68", "wheel"]
build-backend = "setuptools.build_meta"

[project]
name = "myproject"
version = "0.1.0"
description = "My project"
readme = "README.md"
requires-python = ">=3.11"
license = {text = "MIT"}
authors = [{name = "Ivan", email = "ivan@example.com"}]
dependencies = [
    "requests>=2.31",
    "pydantic>=2.6",
]

[project.optional-dependencies]
dev = [
    "pytest>=8.0",
    "ruff>=0.4",
    "mypy>=1.10",
]

[project.scripts]
myproject = "myproject.cli:main"

[tool.ruff]
line-length = 100

[tool.mypy]
strict = true

[tool.pytest.ini_options]
testpaths = ["tests"]

Плюсы:
Единое место для метаданных и настроек инструментов.
Декларативный.
Совместим с pip, uv, poetry, hatch, pdm.
Lock-файлы (uv.lock, poetry.lock) дополняют.

Минусы:
Сложнее для новичка.
Разные бэкенды (setuptools, hatchling, flit, poetry-core).

Когда использовать: всегда в новых проектах.

5.4. Сравнение
ФайлНазначениеСтатусsetup.pyсборка/установкаlegacyrequirements.txtсписок зависимостейактуален для простых случаевpyproject.tomlвсё: метаданные, зависимости, настройкисовременный стандартuv.lock / poetry.lockточные версии всех зависимостейобязателен для воспроизводимости

5.5. Что проверить руками

mkdir /tmp/pyproj && cd /tmp/pyproj
cat > pyproject.toml <<'EOF'
[project]
name = "demo"
version = "0.1.0"
requires-python = ">=3.11"
dependencies = ["requests"]
EOF

python3 -m pip install -e .
python3 -c "import demo"  # если есть пакет

6. Git, .gitignore для Python
6.1. Базовый workflow

git init
git add .
git commit -m "Initial commit"
git status
git log --oneline
git diff
git branch feature/x
git checkout feature/x
git merge feature/x

6.2. .gitignore для Python

Что не коммитить:

# Виртуальные окружения
.venv/
venv/
env/
ENV/

# Python
__pycache__/
*.py[cod]
*$py.class
*.so

# Сборка
build/
dist/
*.egg-info/
*.egg
.eggs/

# Тесты и кэш
.pytest_cache/
.coverage
.coverage.*
htmlcov/
.tox/
.nox/
.hypothesis/

# Типы
.mypy_cache/
.pytype/
.ruff_cache/

# Jupyter
.ipynb_checkpoints/

# IDE
.vscode/
.idea/
*.swp
*.swo

# OS
.DS_Store
Thumbs.db

# Env-файлы (секреты!)
.env
.env.*
!.env.example

# Логи
*.log

Важно:
Секреты (.env, токены, ключи) — никогда в git.
Вместо .env коммитить .env.example с пустыми значениями.
.venv — всегда игнорировать.
__pycache__ — всегда игнорировать.

6.3. Проверка перед коммитом

git status --ignored
git check-ignore -v .venv/
git ls-files | grep -E "(__pycache__|\.venv|\.env$)"Если что-то лишнее попало — удалить из индекса:git rm -r --cached .venv
git rm --cached .env6.4. Полезные алиасыgit config --global alias.st status
git config --global alias.co checkout
git config --global alias.br branch
git config --global alias.lg "log --oneline --graph --all"6.5. Что проверить рукамиmkdir /tmp/gittest && cd /tmp/gittest
git init
python3 -m venv .venv
source .venv/bin/activate
python3 -c "print('hello')" > main.py
python3 main.py
python3 -c "import main" 2>/dev/null; ls __pycache__ 2>/dev/null

# Создать .gitignore, проверить
echo ".venv/" > .gitignore
echo "__pycache__/" >> .gitignore
git add .
git status

7. Редактор: VS Code + Python extension или PyCharm
7.1. VS Code

Плюсы:
Бесплатный, лёгкий, кроссплатформенный.
Огромная экосистема расширений.
Отличная поддержка Python.
Работает с любым venv.

Обязательные расширения:
Python (Microsoft) — IntelliSense, отладка, Jupyter, тесты.
Pylance — быстрый language server, типы.
Ruff — линтер + форматтер.
Mypy Type Checker — проверка типов.
Even Better TOML — для pyproject.toml.
GitLens — история git.
Error Lens — ошибки прямо в строке.
Настройки (.vscode/settings.json):
{
  "python.defaultInterpreterPath": ".venv/bin/python",
  "python.analysis.typeCheckingMode": "strict",
  "editor.formatOnSave": true,
  "[python]": {
    "editor.defaultFormatter": "charliermarsh.ruff",
    "editor.codeActionsOnSave": {
      "source.organizeImports": "explicit"
    }
  },
  "python.testing.pytestEnabled": true,
  "python.testing.pytestArgs": ["tests"]
}Отладка (launch.json):{
  "version": "0.2.0",
  "configurations": [
    {
      "name": "Python: Current File",
      "type": "debugpy",
      "request": "launch",
      "program": "${file}",
      "console": "integratedTerminal",
      "justMyCode": true
    }
  ]
}
Плюсы VS Code:
Лёгкий.
Гибкий.
Отличная поддержка Python.
Бесплатный.

Минусы:
Нужно настраивать расширения.
Иногда конфликты между расширениями.

7.2. PyCharm
Плюсы:
Мощный из коробки.
Отличный рефакторинг.
Встроенная поддержка venv, poetry, uv.
Профессиональный debugging.
Database tools (в Professional).
Scientific mode (Professional).
Версии:
Community — бесплатная, для чистого Python.
Professional — платная, для web (Django, Flask, FastAPI), БД, science.

Настройки:
File → Settings → Project → Python Interpreter → выбрать .venv.
Settings → Tools → Ruff / Black.
Settings → Editor → Inspections → MyPy.
Плюсы PyCharm:
«Всё из коробки».
Мощный рефакторинг.
Отличная навигация.

Минусы:
Тяжёлый (много RAM).
Медленный старт.
Professional платный.

7.3. Что выбратьКритерийVS CodePyCharmВеслёгкийтяжёлый
Настройкавручнуюиз коробкиРефакторингхорошоотличноОтладкахорошоотличноЦенабесплатноCommunity бесплатно, Pro платноДля новичкахорошоотличноДля профиотличноотлично

Рекомендация:
Новичку — VS Code (лёгкий, бесплатный, учит настраивать).
Профи — PyCharm Professional (если есть бюджет) или VS Code.

7.4. Что проверить руками

VS Code:

code --install-extension ms-python.python
code --install-extension ms-python.vscode-pylance
code --install-extension charliermarsh.ruff
code --install-extension ms-python.mypy-type-checker

Открыть проект, выбрать интерпретатор .venv, поставить breakpoint, запустить F5.
PyCharm:
File → Settings → Project → Python Interpreter → Add → Existing → .venv/bin/python.
Run → Edit Configurations → Python → указать скрипт.

8. Terminal: bash/zsh, базовые команды
8.1. Что такое shell

Shell — программа, которая читает команды и запускает процессы.
bash — Bourne Again Shell, стандарт на Linux.
zsh — Z Shell, стандарт на macOS, мощнее bash.
fish — дружелюбный, но не POSIX.
sh — минимальный POSIX shell.

Проверка:

echo $SHELL
echo $0

8.2. Базовые команды

Навигация:
pwd                 # текущая директория
cd /path            # перейти
cd ..               # на уровень выше
cd ~                # в домашнюю
cd -                # в предыдущую
ls                  # список файлов
ls -la              # с деталями и скрытыми
tree                # дерево (если установлено)Файлы:touch file.txt
mkdir -p a/b/c
cp file.txt copy.txt
cp -r dir/ dir2/
mv file.txt new.txt
rm file.txt
rm -rf dir/
cat file.txt
less file.txt
head -n 20 file.txt
tail -n 20 file.txt
tail -f log.txtПоиск:grep "pattern" file.txt
grep -r "pattern" .
find . -name "*.py"
find . -type f -size +1M
which python3
whereis python3Права:chmod +x script.sh
chmod 644 file.txt
chown user:group file.txtПроцессы:ps aux
top
htop
kill PID
kill -9 PID
pkill -f "python"Сеть:curl https://example.com
wget https://example.com/file
ping google.com
nc -zv localhost 8080
ss -tulnПайпы и редиректы:ls | grep ".py"
cat file.txt | wc -l
echo "hello" > file.txt
echo "world" >> file.txt
command 2>&1 | tee log.txtПеременные:export MY_VAR=value
echo $MY_VAR
env
unset MY_VAR

История:

history
!123                # выполнить команду №123
!!                  # предыдущая
Ctrl+R              # поиск по историиСпециальные:Ctrl+C              # прервать
Ctrl+D              # EOF / выход
Ctrl+Z              # приостановить (fg, bg)
Ctrl+L              # очистить экран
Tab                 # автодополнение8.3. Полезные настройкиbash (~/.bashrc):alias ll='ls -la'
alias gs='git status'
alias gp='git pull'
alias venv='python3 -m venv .venv && source .venv/bin/activate'
export EDITOR=vim
export PYTHONDONTWRITEBYTECODE=1zsh (~/.zshrc):alias ll='ls -la'
alias gs='git status'
export EDITOR=vim

Применить:

source ~/.bashrc
# или
source ~/.zshrc

8.4. Python-специфичные команды

# Активация venv
source .venv/bin/activate

# Проверка, что python из venv
which python
python -c "import sys; print(sys.prefix)"

# Запуск модуля
python -m pytest
python -m mypy src/
python -m ruff check .

# Отладка
python -m pdb script.py
python -X dev script.py     # dev mode
python -X faulthandler script.py8.5. Что проверить руками# Создать структуру
mkdir -p /tmp/termtest/{src,tests}
cd /tmp/termtest
touch README.md src/main.py tests/test_main.py

# Навигация
pwd
ls -la
tree

# Поиск
find . -name "*.py"
grep -r "main" .

# Пайпы
find . -name "*.py" | wc -l
find . -name "*.py" | xargs wc -l

# Процессы
ps aux | grep python

# История
history | tail -209.

Итоговая шпаргалка: что нужно уметь после Лекции 0

Понимать:
Как Python превращает .py в байт-код и выполняет его.
Зачем нужен REPL и когда его использовать.
Что такое __pycache__ и почему его не коммитят.
Почему Python 2 мёртв и какие версии актуальны.
Разницу между python, python3, pip, pipx, uv, poetry.
Зачем изоляция зависимостей и как работает venv.
Разницу между setup.py, requirements.txt, pyproject.toml.
Что не коммитить в git (.venv, __pycache__, .env).
Как настроить VS Code или PyCharm под Python.
Базовые команды терминала.

Уметь:
Создать venv, активировать, установить пакет, деактивировать.
Создать pyproject.toml с зависимостями.
Установить проект в editable-режиме (pip install -e .).
Настроить .gitignore для Python.
Сделать первый коммит.
Запустить Python-скрипт и REPL.
Посмотреть байт-код через dis.
Найти файл через find, строку через grep.
Работать с пайпами и редиректами.

Проверить руками (минимальный чек-лист):

mkdir /tmp/checkpoint0 && cd /tmp/checkpoint0
python3 -m venv .venv
source .venv/bin/activate
python3 --version
pip install --upgrade pip
pip install requests
python3 -c "import requests; print(requests.__version__)"
deactivate

git init
cat > .gitignore <<'EOF'
.venv/
__pycache__/
*.pyc
.env
EOF

cat > main.py <<'EOF'
def main() -> None:
    print("Hello, World!")

if __name__ == "__main__":
    main()
EOF

python3 main.py
git add .
git commit -m "Initial commit"
git log --oneline10. 

Что читать дальше

PEP 8 — стиль кода.
PEP 484 — типы.
PEP 517/518/621 — pyproject.toml.
PEP 405 — venv.
Python docs: Tutorial — https://docs.python.org/3/tutorial/Real Python — практические статьи.
Python Packaging User Guide — https://packaging.python.org/
Этот конспект — фундамент. 
Дальше каждый этап будет опираться на эти понятия: venv, pyproject.toml, git, ruff, mypy. 
Если что-то непонятно — вернись к соответствующему разделу и проверь руками.

ПРАКТИКА:

Задача 0.1. hello с правильным окружением
L1 — работает:
mkdir hello && cd hello
python3 -m venv .venv
source .venv/bin/activate
pip install --upgrade pip# hello.py
def main() -> None:
    print("Hello, World!")

if __name__ == "__main__":
    main()python hello.py

L2 — инженерный:pyproject.toml:
[project]
name = "hello"
version = "0.1.0"
requires-python = ">=3.10"

[tool.ruff]
line-length = 100

[tool.mypy]
strict = true.gitignore для Python.git init, первый коммит.

L3 — production-minded:
CI (GitHub Actions): запуск ruff, mypy, pytest на push.
README с инструкцией.Что переиспользуем дальше: структуру проекта, pyproject.toml, CI.

Типичные ошибки:
Ставить пакеты глобально, а не в venv
.Коммитить .venv.
Использовать python вместо python3 на Linux.